% ==========================================
% BAB I PENDAHULUAN
% ==========================================
\cleardoublepage
\chapter{PENDAHULUAN}
\label{chap:I}

\section{Latar Belakang}
\label{sec:I.1}

Kementerian Pertahanan Republik Indonesia (Kemhan RI) merupakan kementerian yang menyelenggarakan urusan pemerintahan di bidang pertahanan \cite{Permenhan13Tahun2026}. Berdasarkan Undang-Undang Nomor 3 Tahun 2002 tentang Pertahanan Negara, Menteri Pertahanan membantu Presiden dalam merumuskan kebijakan umum pertahanan negara dan menetapkan kebijakan penyelenggaraan pertahanan negara. Menteri juga menetapkan kebijakan penganggaran, pengadaan, perekrutan, pengelolaan sumber daya nasional, serta pembinaan teknologi dan industri pertahanan \cite{UU3Tahun2002}. Untuk menjalankan mandat tersebut, Kemhan RI membutuhkan data dan informasi pertahanan yang lengkap, akurat, dan mutakhir, seperti data personel, alutsista, logistik, anggaran, dan informasi ancaman.

Kebutuhan tersebut semakin mendesak karena kondisi saat ini ditandai oleh ancaman yang bersifat multidimensi, mulai dari ancaman militer konvensional hingga ancaman siber dan disrupsi informasi. Keunggulan pertahanan tidak lagi semata ditentukan oleh jumlah alutsista dan personel, tetapi juga oleh kemampuan mengintegrasikan sistem dan pengambilan keputusan secara cepat berbasis data \cite{USDoD2019,Koch2015,USDoD2022}. Di tingkat nasional, Pemerintah juga mendorong setiap kementerian, termasuk Kemhan RI, untuk membangun ekosistem digital yang terintegrasi melalui Sistem Pemerintahan Berbasis Elektronik (SPBE) \cite{Perpres95Tahun2018}, percepatan transformasi digital \cite{Perpres82Tahun2023}, dan RPJPN 2025–2045 \cite{UU59Tahun2024}. Dengan demikian, Kemhan RI memiliki dua kebutuhan sekaligus. Sebagai kementerian, Kemhan RI wajib menyusun Arsitektur SPBE yang berorientasi pada layanan administrasi pemerintahan. Sebagai institusi pertahanan, Kemhan RI juga perlu menjangkau kapabilitas operasional pertahanan, terutama keterpaduan komando dan kendali lintas matra TNI.

Dalam praktiknya, transformasi digital di lingkungan Kemhan RI masih menghadapi permasalahan utama berupa sistem yang berkembang secara terpisah (\textit{silo}) di setiap satuan kerja dan matra. \textit{Silo} tersebut terjadi pada beberapa lapisan Kemhan RI, yaitu:
\begin{enumerate}[label=\arabic*.]
\item Lapisan bisnis dan tata kelola. Pengelolaan data, sistem informasi, dan infrastruktur teknologi informasi dan komunikasi telah ditetapkan sebagai tugas Pusat Data dan Informasi \cite{Permenhan30Tahun2025}. Namun, mandat unit tersebut bersifat menjalankan sistem yang sudah ada (\textit{run-the-engine}), dan belum ada fungsi yang mengorkestrasi perubahan lintas organisasi (\textit{change-the-business}). Akibatnya, setiap satuan kerja merencanakan dan mengembangkan sistemnya sendiri.
\item Lapisan data dan aplikasi. Inventarisasi aplikasi Kemhan RI tahun 2025 mencatat 247 aplikasi, terdiri atas 176 aplikasi umum dan 71 aplikasi khusus. Aplikasi khusus dikembangkan per satuan kerja, dan sebagian memiliki fungsi serupa, antara lain pada kepegawaian, pembelajaran daring, dan persuratan. Meskipun kebijakan Satu Data Pertahanan telah ditetapkan \cite{Permenhan1Tahun2023}, data masih terkotak per aplikasi karena belum tersedia kamus data, klasifikasi data, dan kanal berbagi pakai.
\item Lapisan teknologi. Seluruh layanan bertumpu pada satu \textit{Data Center} tanpa \textit{Disaster Recovery Center}, sehingga menjadi titik tunggal kegagalan (\textit{single point of failure}). Lapisan keamanan juga masih bersifat parsial, padahal ancaman siber terhadap infrastruktur dan data pertahanan terus meningkat \cite{Permenhan82Tahun2014}.
\item Lapisan operasional. Komando dan kendali masih dikembangkan per matra. Kondisi ini menyulitkan keterpaduan antara unsur ISR, komando-kendali, dan unsur penindak dalam operasi gabungan.
\end{enumerate}

Kondisi \textit{silo} pada keempat lapisan tersebut mengakibatkan beberapa hal, yaitu:
\begin{enumerate}[label=\alph*.]
\item investasi digital berulang dan penggunaan sumber daya yang tidak efisien;
\item data sulit dibagi pakai dan diintegrasikan untuk analisis lintas domain;
\item pengambilan keputusan strategis berjalan lambat;
\item interoperabilitas rendah, baik antarsatuan kerja maupun antarmatra; dan
\item layanan dan data pertahanan rentan terhadap gangguan dan serangan siber.
\end{enumerate}

Permasalahan tersebut terjadi karena Kemhan RI belum memiliki arsitektur yang menjadi acuan dalam membangun sistem secara menyeluruh, mulai dari proses bisnis, data, aplikasi, teknologi, hingga kapabilitas operasional. Indeks SPBE Kemhan RI yang meningkat dari 2,62 pada tahun 2022 menjadi 3,43 pada tahun 2023 dengan predikat Baik \cite{Pusdatin2024}, menunjukkan bahwa tata kelola SPBE telah berjalan. Namun, indeks tersebut tidak dirancang untuk mengukur keterpaduan kapabilitas operasional pertahanan. Oleh karena itu, diperlukan solusi yang mampu menyatukan sisi administratif dan sisi operasional dalam satu arah transformasi.

Salah satu solusi untuk mengatasi permasalahan \textit{silo} adalah pengembangan \textit{Enterprise Architecture} (EA). EA menyelaraskan strategi organisasi dengan proses bisnis, data, aplikasi, dan teknologi secara terstruktur, sehingga setiap pengembangan sistem mengacu pada satu rancangan yang sama. Berbagai penelitian menunjukkan peran EA sebagai instrumen transformasi digital \cite{Syamfithriani2024}, termasuk dalam transformasi organisasi militer \cite{Mattila2022}. Pendekatan ini juga sejalan dengan kewajiban penyusunan Arsitektur SPBE bagi instansi pemerintah \cite{Perpres95Tahun2018}. Agar EA dapat dikembangkan secara terarah, organisasi perlu mengadopsi sebuah \textit{framework} EA yang menjadi panduan pengembangannya.

Beberapa \textit{framework} EA telah dikembangkan khusus untuk lingkungan pertahanan, yaitu DoDAF, MODAF, dan NAF \cite{DoDAF2010v202,Weilkiens2022,NAFv41}. Ketiganya berakar pada C4ISR \textit{Architecture Framework} \cite{C4ISRAWG1998}, sehingga kuat dalam memuat aspek komando, kendali, dan interoperabilitas. Namun, DoDAF dan MODAF pada dasarnya merupakan kerangka deskripsi yang menetapkan apa yang harus didokumentasikan, bukan bagaimana arsitektur dibangun secara bertahap. NAF telah dilengkapi metodologi sejak versi 4, tetapi dirancang untuk interoperabilitas multinasional dan mengikuti kebijakan NATO \cite{NAFv41}.

Di sisi lain, TOGAF merupakan salah satu \textit{framework} EA dengan adopsi terluas dan tersedia secara terbuka untuk penggunaan akademik dan nonkomersial. TOGAF memiliki \textit{Architecture Development Method} (ADM), yaitu metode pengembangan arsitektur yang bertahap dan iteratif \cite{TOGAF2022}. Salah satu karakteristik TOGAF adalah sifatnya yang generik dan fleksibel, sehingga TOGAF mewajibkan penyesuaian (\textit{tailoring}) sejak Fase \textit{Preliminary}. Namun, sifat generik tersebut membuat TOGAF tidak menyediakan muatan pertahanan, seperti domain kapabilitas komando dan kendali, penautan arsitektur dengan doktrin pertahanan, dan pola integrasi lintas matra. Keamanan siber pun masih diperlakukan sebagai \textit{concern} lintas domain \cite{TOGAFRiskSecurity}. Tanpa acuan muatan pertahanan, penyesuaian TOGAF bergantung pada improvisasi setiap arsitek, sehingga arsitektur yang dihasilkan cenderung hanya mencakup sisi administratif. Dengan demikian, perancangan EA menggunakan TOGAF saja belum cukup menjawab kebutuhan Kemhan RI, dan masih dibutuhkan sebuah model rujukan pertahanan yang dapat ditanamkan ke dalam TOGAF.

Model rujukan yang sesuai untuk kebutuhan tersebut adalah C6ISR, yaitu \textit{Command, Control, Communications, Computers, Cyber-Defense, Combat Systems, Intelligence, Surveillance, and Reconnaissance}. C6ISR merupakan perkembangan bertahap dari konsep komando dan kendali (C2) melalui C4ISR dan C5ISR \cite{Petrovski2023,Todorovic2025}. Dibandingkan C4ISR, C6ISR menambahkan unsur \textit{cyber-defense} sebagai domain operasi \cite{NATO2016Warsaw}, serta unsur \textit{combat systems} yang melengkapi siklus operasi dari pengindraan dan pengambilan keputusan hingga penindakan beserta umpan baliknya.

Dengan adanya TOGAF sebagai \textit{framework} yang unggul dalam metode pengembangan arsitektur, serta C6ISR sebagai model rujukan kapabilitas pertahanan, penelitian ini memodifikasi \textit{framework} TOGAF untuk Kemhan RI yang berbasis pada C6ISR. Modifikasi dilakukan dengan mempertahankan siklus ADM, menambahkan Fase \textit{C6ISR Architecture} sebagai domain arsitektur tersendiri, serta menyesuaikan muatan fase lainnya agar sesuai dengan kebutuhan pertahanan. \textit{Framework} hasil modifikasi tersebut selanjutnya disebut TOGAF Pertahanan Indonesia, kemudian diterapkan pada Kemhan RI untuk menghasilkan arsitektur transformasi digital beserta \textit{roadmap} implementasinya.

\section{Rumusan Masalah}
\label{sec:I.2}

Berdasarkan latar belakang tersebut, permasalahan utama penelitian ini adalah belum tersedianya \textit{framework} EA yang mampu memenuhi kewajiban Arsitektur pemerintahan berbasis elektronik sekaligus menjangkau kapabilitas operasional pertahanan Kemhan RI. Permasalahan tersebut dirumuskan dalam tiga pertanyaan penelitian (\textit{research question}, RQ) berikut:
\begin{enumerate}[label=RQ\arabic*:, leftmargin=*]
\item Bagaimana memodifikasi TOGAF ADM berbasis C6ISR agar sesuai dengan kebutuhan arsitektur enterprise Kemhan RI?
\item Bagaimana penerapan \textit{framework} hasil modifikasi pada Kemhan RI dalam menghasilkan arsitektur transformasi digital beserta \textit{roadmap} implementasinya?
\item Sejauh mana \textit{framework} hasil modifikasi dan arsitektur yang dihasilkannya layak diterapkan, serta apa kontribusinya dibandingkan TOGAF ADM standar?
\end{enumerate}

\section{Tujuan Penelitian}
\label{sec:I.3}

Tujuan penelitian ini adalah sebagai berikut:
\begin{enumerate}[label=\arabic*.]
\item Merancang modifikasi TOGAF ADM berbasis C6ISR yang sesuai dengan kebutuhan Kemhan RI, selanjutnya disebut TOGAF Pertahanan Indonesia.
\item Menerapkan TOGAF Pertahanan Indonesia pada Kemhan RI untuk menghasilkan arsitektur transformasi digital beserta \textit{roadmap} implementasi 2026–2030.
\item Mengevaluasi kelayakan \textit{framework} dan arsitektur yang dihasilkan, serta kontribusinya dibandingkan TOGAF ADM standar.
\end{enumerate}

\section{Manfaat Penelitian}
\label{sec:I.4}

Penelitian ini diharapkan memberikan manfaat sebagai berikut:
\begin{enumerate}[label=\arabic*.]
\item Bagi pengembangan ilmu, penelitian ini menghasilkan modifikasi TOGAF berbasis C6ISR yang terdokumentasi dan dapat diulang, sehingga dapat menjadi rujukan pengembangan EA di sektor pertahanan.
\item Bagi Kemhan RI, penelitian ini menyediakan \textit{framework} serta arsitektur transformasi digital beserta \textit{roadmap} implementasinya sebagai acuan dalam mengintegrasikan sistem, data, serta komando dan kendali.
\end{enumerate}

\section{Ruang Lingkup Penelitian}
\label{sec:I.5}

Ruang lingkup penelitian ini adalah sebagai berikut:
\begin{enumerate}[label=\arabic*.]
\item Penelitian ini memodifikasi TOGAF ADM agar sesuai dengan kebutuhan pertahanan. Modifikasi ini mencakup penyesuaian masukan, langkah kerja, dan keluaran pada setiap fase serta penambahan fase C6ISR.
\item Penerapan dilakukan dari Fase \textit{Preliminary} hingga Fase \textit{Migration Planning}. Fase \textit{Implementation Governance} dan \textit{Architecture Change Management} hanya dirancang dan tidak diterapkan.
\item Penelitian tidak mencakup aspek rahasia militer dan operasi tempur, kecuali pada Fase \textit{C6ISR Architecture} yang memodelkan keterhubungan komando dan kendali dengan sistem tempur secara konseptual berdasarkan doktrin, kebijakan, dan sumber terbuka. Fase tersebut dirancang dari sudut pandang Kemhan RI sebagai penyelenggara kebijakan dan standar, sedangkan komando dan pengendalian operasional tetap berada pada Mabes TNI.
\item Evaluasi bersifat \textit{ex-ante}, yaitu menilai rancangan sebelum diimplementasikan, sehingga dampak implementasi arsitektur tidak diukur.
\end{enumerate}

\section{Hasil / Keluaran yang Diharapkan}
\label{sec:I.6}

Berdasarkan tujuan penelitian, hasil yang diharapkan adalah sebagai berikut:
\begin{enumerate}[label=\arabic*.]
\item TOGAF Pertahanan Indonesia, yaitu \textit{framework} hasil modifikasi TOGAF ADM yang memuat Fase \textit{C6ISR Architecture} sebagai fase baru dan penyesuaian isi lima fase lainnya agar sesuai dengan kebutuhan pertahanan Indonesia.
\item Arsitektur transformasi digital Kemhan RI yang meliputi arsitektur bisnis, data, aplikasi, teknologi, dan C6ISR, beserta paket pekerjaan, \textit{roadmap}, dan rencana migrasi 2026–2030.
\item Hasil evaluasi kelayakan \textit{framework} dan arsitektur berdasarkan \textit{Enterprise Architecture Scorecard}, penilaian pakar, dan perbandingan dengan TOGAF ADM standar, dengan target kategori Baik pada \textit{Enterprise Architecture Scorecard} dan kategori Layak pada penilaian pakar.
\end{enumerate}

\section{Metodologi Penelitian}
\label{sec:I.7}

\subsection{Pendekatan Penelitian}
\label{sec:I.7.1}

Penelitian ini menggunakan pendekatan \textit{Design Science Research Methodology} (DSRM). Pendekatan ini bertujuan menghasilkan artefak yang dapat digunakan untuk menyelesaikan permasalahan nyata dan dapat diuji kegunaannya. Artefak tersebut dapat berupa konstruk, model, metode, atau instansiasi \cite{Hevner2004}. Dalam penelitian ini, artefak utama berupa metode, yaitu TOGAF Pertahanan Indonesia, sedangkan arsitektur transformasi digital Kemhan RI merupakan instansiasi yang digunakan untuk mendemonstrasikan metode tersebut. DSRM dipilih karena \textit{framework} hasil modifikasi perlu dirancang berdasarkan kebutuhan, didemonstrasikan pada kondisi nyata, dan dievaluasi oleh pakar.

\subsection{Tahapan Penelitian}
\label{sec:I.7.2}

Penelitian ini mengikuti enam tahapan DSRM menurut \textcite{Peffers2007}, yaitu sebagai berikut.

\subsubsection{Tahap 1 – Problem Identification and Motivation}
\label{sec:I.7.2.1}

Tahap ini mengidentifikasi permasalahan dan urgensi penelitian melalui analisis dokumen dan regulasi, wawancara dengan unit kerja Kemhan RI, serta kajian penelitian terdahulu. Hasilnya adalah permasalahan \textit{silo} pada lapisan bisnis, data dan aplikasi, teknologi, serta operasional Kemhan RI, dan belum tersedianya \textit{framework} EA yang menjangkau sisi administratif sekaligus sisi operasional pertahanan.

\subsubsection{Tahap 2 – Defining the Objectives for a Solution}
\label{sec:I.7.2.2}

Tahap ini menetapkan sasaran yang harus dipenuhi oleh \textit{framework} hasil modifikasi. Sasaran diturunkan dari tiga sumber, yaitu karakteristik domain pertahanan, regulasi nasional, dan praktik terbaik internasional, kemudian dirumuskan sebagai kebutuhan \textit{framework}. Secara ringkas, \textit{framework} harus mempertahankan metode ADM yang bertahap dan iteratif, mengakomodasi C6ISR sebagai domain arsitektur, menanamkan keamanan siber sejak perancangan, mendukung interoperabilitas lintas satuan kerja dan matra, serta tetap selaras dengan SPBE.
\begin{landscape}

\subsubsection{Tahap 3 – Design and Development}
\label{sec:I.7.2.3}

Tahap ini merancang TOGAF Pertahanan Indonesia, TOGAF ADM digunakan sebagai dasar, kemudian ditambahkan Fase \textit{C6ISR Architecture} sebagai fase baru dan disesuaikan muatan lima fase lainnya.Hasilnya disajikan pada Tabel~\ref{tab:I.1}.

\begingroup\small\setstretch{1.0}
\begin{longtable}{@{}>{\raggedright\arraybackslash\hspace{0pt}}p{0.044\linewidth}>{\raggedright\arraybackslash\hspace{0pt}}p{0.170\linewidth}>{\raggedright\arraybackslash\hspace{0pt}}p{0.244\linewidth}>{\raggedright\arraybackslash\hspace{0pt}}p{0.292\linewidth}>{\raggedright\arraybackslash\hspace{0pt}}p{0.176\linewidth}@{}}
\caption{Fase TOGAF Pertahanan Indonesia dalam Lingkup Penelitian}\label{tab:I.1}\\
\toprule
\textbf{No.} & \textbf{Tahapan TOGAF} & \textbf{Tujuan Utama} & \textbf{Aktifitas Kunci} & \textbf{Hasil Akhir}\\
\midrule
\endfirsthead
\multicolumn{5}{@{}l}{\small\itshape \tablename\ \thetable{} (lanjutan)}\\
\toprule
\textbf{No.} & \textbf{Tahapan TOGAF} & \textbf{Tujuan Utama} & \textbf{Aktifitas Kunci} & \textbf{Hasil Akhir}\\
\midrule
\endhead
\bottomrule
\endlastfoot
\centering\arraybackslash 1 & \centering\arraybackslash 2 & \centering\arraybackslash 3 & \centering\arraybackslash 4 & \centering\arraybackslash 5\\
\centering\arraybackslash 1 & \textit{Preliminary Phase} (Δ) & Menetapkan fondasi pengembangan arsitektur, yaitu ruang lingkup, prinsip, pemangku kepentingan, dan tata kelola arsitektur. & \hangindent=1.1em\hangafter=1 \makebox[1.1em][l]{--}Menautkan doktrin, strategi pertahanan, dan regulasi nasional sebagai masukan.\par \hangindent=1.1em\hangafter=1 \makebox[1.1em][l]{--}Menetapkan sembilan prinsip arsitektur Pertahanan Indonesia.\par \hangindent=1.1em\hangafter=1 \makebox[1.1em][l]{--}Menetapkan ruang lingkup arsitektur dan mengidentifikasi pemangku kepentingan.\par \hangindent=1.1em\hangafter=1 \makebox[1.1em][l]{--}Menetapkan tata kelola dan repositori arsitektur. & Prinsip arsitektur, ruang lingkup, peta pemangku kepentingan, dan model tata kelola arsitektur.\\
\centering\arraybackslash 2 & \textit{Architecture Vision} (Δ) & Merumuskan visi dan tujuan arsitektur yang berorientasi pada pertahanan cerdas (\textit{smart defence}). & \hangindent=1.1em\hangafter=1 \makebox[1.1em][l]{--}Menyusun visi, misi, dan tujuan arsitektur.\par \hangindent=1.1em\hangafter=1 \makebox[1.1em][l]{--}Mengidentifikasi pendorong dan kendala transformasi.\par \hangindent=1.1em\hangafter=1 \makebox[1.1em][l]{--}Menilai kapabilitas, kesiapan, dan risiko transformasi.\par \hangindent=1.1em\hangafter=1 \makebox[1.1em][l]{--}Menyusun konsep solusi dan \textit{Statement of Architecture Work}. & Visi arsitektur dan \textit{Statement of Architecture Work}.\\
\centering\arraybackslash 3 & \textit{Business Architecture} (Δ) & Merancang struktur organisasi dan kapabilitas bisnis yang mendukung komando-kendali serta kolaborasi lintas satuan. & \hangindent=1.1em\hangafter=1 \makebox[1.1em][l]{--}Menganalisis struktur organisasi dan tata kelola.\par \hangindent=1.1em\hangafter=1 \makebox[1.1em][l]{--}Menyusun \textit{value stream} dan peta kapabilitas.\par \hangindent=1.1em\hangafter=1 \makebox[1.1em][l]{--}Merancang struktur organisasi dan peta kapabilitas target\par \hangindent=1.1em\hangafter=1 \makebox[1.1em][l]{--}Menganalisis kesenjangan dan menyusun \textit{roadmap}. & Peta kapabilitas dan arsitektur bisnis target.\\
\centering\arraybackslash 4 & \textit{Information Systems Architecture} (Δ) & Merancang arsitektur data dan aplikasi yang terintegrasi dan aman. & \hangindent=1.1em\hangafter=1 \makebox[1.1em][l]{--}Menyusun model entitas, klasifikasi kerahasiaan data, kamus data, serta kepemilikan data.\par \hangindent=1.1em\hangafter=1 \makebox[1.1em][l]{--}Menganalisis portofolio dan integrasi aplikasi serta melakukan rasionalisasi.\par \hangindent=1.1em\hangafter=1 \makebox[1.1em][l]{--}Merancang arsitektur data dan aplikasi target beserta interoperabilitasnya.\par \hangindent=1.1em\hangafter=1 \makebox[1.1em][l]{--}Menganalisis kesenjangan dan menyusun \textit{roadmap}. & Arsitektur data dan arsitektur aplikasi target.\\
\centering\arraybackslash 5 & \textit{Technology Architecture} (Δ) & Merancang infrastruktur teknologi dan fondasi keamanan siber bagi ekosistem digital pertahanan. & \hangindent=1.1em\hangafter=1 \makebox[1.1em][l]{--}Menganalisis arsitektur teknologi \textit{baseline}.\par \hangindent=1.1em\hangafter=1 \makebox[1.1em][l]{--}Menetapkan prinsip dan standar teknologi.\par \hangindent=1.1em\hangafter=1 \makebox[1.1em][l]{--}Merancang arsitektur teknologi target dan lapisan keamanan.\par \hangindent=1.1em\hangafter=1 \makebox[1.1em][l]{--}Menganalisis kesenjangan dan menyusun \textit{roadmap}. & Arsitektur teknologi target dan arsitektur keamanan.\\
\centering\arraybackslash 6 & \textit{C6ISR Architecture} (+) & Merancang kapabilitas C6ISR yang mengintegrasikan ISR, komando-kendali, dan unsur penindak untuk mendukung operasi lintas matra. & \hangindent=1.1em\hangafter=1 \makebox[1.1em][l]{--}Mengidentifikasi kebutuhan operasional dan arsitektur C6ISR \textit{baseline}.\par \hangindent=1.1em\hangafter=1 \makebox[1.1em][l]{--}Menyusun peta kapabilitas C6ISR.\par \hangindent=1.1em\hangafter=1 \makebox[1.1em][l]{--}Merancang arsitektur target per lapisan (\textit{Command \& Control}, \textit{Computers}, \textit{Communications}, \textit{Cyber Defense}, \textit{Combat Systems}, dan ISR), termasuk skema komunikasi PACE.\par \hangindent=1.1em\hangafter=1 \makebox[1.1em][l]{--}Menganalisis kesenjangan dan menyusun \textit{roadmap}. & Arsitektur C6ISR target.\\
\centering\arraybackslash 7 & \textit{Opportunities \& Solutions} & Mengonsolidasikan kesenjangan seluruh domain menjadi paket pekerjaan yang siap diimplementasikan. & \hangindent=1.1em\hangafter=1 \makebox[1.1em][l]{--}Mengidentifikasi atribut perubahan dan kendala implementasi.\par \hangindent=1.1em\hangafter=1 \makebox[1.1em][l]{--}Mengonsolidasikan kesenjangan lintas domain.\par \hangindent=1.1em\hangafter=1 \makebox[1.1em][l]{--}Merumuskan paket pekerjaan dan ketergantungannya.\par \hangindent=1.1em\hangafter=1 \makebox[1.1em][l]{--}Menilai kesiapan dan risiko serta menyusun arsitektur transisi. & Paket pekerjaan, arsitektur transisi, dan \textit{roadmap} arsitektur.\\
\centering\arraybackslash 8 & \textit{Migration Planning} & Menyusun rencana implementasi dan migrasi transformasi digital secara bertahap. & \hangindent=1.1em\hangafter=1 \makebox[1.1em][l]{--}Menetapkan model tata kelola implementasi.\par \hangindent=1.1em\hangafter=1 \makebox[1.1em][l]{--}Menilai nilai manfaat dan kebutuhan sumber daya.\par \hangindent=1.1em\hangafter=1 \makebox[1.1em][l]{--}Menetapkan prioritas paket pekerjaan.\par \hangindent=1.1em\hangafter=1 \makebox[1.1em][l]{--}Menyusun rencana implementasi dan migrasi 2026–2030 beserta indikator keberhasilan. & \textit{Roadmap} implementasi dan rencana migrasi 2026–2030.\\
\end{longtable}
\endgroup

Keterangan: Δ = muatan fase disesuaikan; + = fase baru; tanpa tanda = mengikuti TOGAF ADM standar.

Tabel~\ref{tab:I.1} menunjukkan bahwa Fase \textit{C6ISR Architecture} merupakan fase baru, sedangkan muatan pertahanan pada Fase \textit{Preliminary} hingga \textit{Technology Architecture} ditambahkan melalui penyesuaian. Fase \textit{Opportunities \& Solutions} dan \textit{Migration Planning} mengikuti TOGAF ADM standar dengan tambahan masukan dari Fase \textit{C6ISR Architecture.}
\end{landscape}

\subsubsection{Tahap 4 – Demonstration}
\label{sec:I.7.2.4}

Tahap ini menerapkan TOGAF Pertahanan Indonesia pada Kemhan RI dari Fase \textit{Preliminary} hingga Fase \textit{Migration Planning}.

\subsubsection{Tahap 5 – Evaluation}
\label{sec:I.7.2.5}

Tahap ini menilai kelayakan \textit{framework} dan arsitektur yang dihasilkan melalui tiga pendekatan berikut:
\begin{enumerate}[label=\arabic*.]
\item Evaluasi teoretis menggunakan \textit{Enterprise Architecture Scorecard} yang diisi oleh pakar akademisi untuk menilai kelengkapan dokumentasi arsitektur.
\item Evaluasi praktis melalui \textit{expert judgement} oleh pakar akademisi, pakar praktisi Kemhan RI, dan pakar mitra, menggunakan FGD dan kuesioner skala Likert 1–5 yang dianalisis dengan rata-rata skor dan indeks Aiken’s V.
\item Evaluasi komparatif yang membandingkan hasil TOGAF Pertahanan Indonesia dengan TOGAF ADM standar untuk menunjukkan kontribusi komponen hasil modifikasi.
\end{enumerate}

\subsubsection{Tahap 6 – Communication}
\label{sec:I.7.2.6}

Tahap ini mengomunikasikan hasil penelitian melalui naskah tesis, dokumen kajian ringkas bagi pakar dan pemangku kepentingan Kemhan RI, serta presentasi kepada Kemhan RI.

\section{Sistematika Penulisan}
\label{sec:I.8}

Tesis ini terdiri atas enam bab dengan uraian sebagai berikut.

\textbf{Bab I Pendahuluan}

Bab ini berisi latar belakang, rumusan masalah, tujuan, manfaat, ruang lingkup, hasil yang diharapkan, metodologi penelitian, dan sistematika penulisan.

\textbf{Bab II Tinjauan Pustaka}

Bab ini berisi penelitian terdahulu, konsep transformasi digital pertahanan, perkembangan komando dan kendali hingga C6ISR, serta perbandingan \textit{framework} EA yang menjadi dasar pemilihan TOGAF.

\textbf{Bab III Analisis dan Desain}

Bab ini berisi analisis kebutuhan \textit{framework} EA Kemhan RI dan perancangan modifikasi TOGAF ADM berbasis C6ISR menjadi TOGAF Pertahanan Indonesia.

\textbf{Bab IV Implementasi}

Bab ini berisi penerapan TOGAF Pertahanan Indonesia pada Kemhan RI dari Fase \textit{Preliminary} hingga Fase \textit{Migration Planning}, termasuk \textit{roadmap} dan rencana migrasi 2026–2030.

\textbf{Bab V Evaluasi}

Bab ini berisi evaluasi \textit{framework} dan arsitektur melalui \textit{Enterprise Architecture Scorecard}, penilaian pakar, dan perbandingan dengan TOGAF ADM standar.

\textbf{Bab VI Kesimpulan}

Bab ini berisi kesimpulan yang menjawab pertanyaan penelitian serta saran bagi Kemhan RI dan penelitian selanjutnya.
